% \begin{figure}[!ht]
%     \centering
%     \includegraphics[width=\linewidth]{figure/validation_accuracy_curve.pdf}
%     \caption{Training efficiency}
%     \label{fig:intro_example}
% \end{figure}

\section{Preliminaries}

\subsection{KL-Based Divergences}
\label{subsec:kl_divergences}


Let $P$ and $Q$ be probability distributions defined over the same sample space $\mathcal{X}$.
The Kullback--Leibler (KL) divergence is a non-symmetric measure of difference between two distributions, quantifying how well $Q$ approximates the reference distribution $P$.
\begin{equation}
\label{eq:kl_general}
\mathrm{KL}(P \,\|\, Q)
= \mathbb{E}_{x \sim P} \left[ \log \frac{P(x)}{Q(x)} \right].
\end{equation}
% The KL divergence is always non-negative and equals zero if and only if $P=Q$.
Due to its asymmetry, $\mathrm{KL}(P \,\|\, Q)$ and $\mathrm{KL}(Q \,\|\, P)$ induce different optimization behaviors, depending on which distribution is used as the reference.

% In this work, we focus on knowledge distillation for autoregressive sequence models, such as large language models, which define a probability distribution over tokens conditioned on a context $\mathbf{c}$.
For auto-regressive sequence models like Large Language Models (LLMs), the probability distribution of a token $x$ is conditioned on a context $\mathbf{c}$, where $\mathbf{c}$ is the sequence generated before $x$. 
% Here, the context $\mathbf{c}$ refers to a prefix of previously generated tokens in the sequence.
% Under this setting, we denote $\pi_\theta(\cdot \mid \mathbf{c})$ as the student model and $\pi_{\tt{te}}(\cdot \mid \mathbf{c})$ as the teacher model.
% Both models define probability distributions over the next token $x$ conditioned on the context $\mathbf{c}$.
In distillation, we have two models, a student model denoted by $\pi_\theta(\cdot \mid \mathbf{c})$, and a teacher model $\pi_{\texttt{te}}(\cdot \mid \mathbf{c})$. We now define forward and reverse KL divergences in terms of these models.

% In the context of knowledge distillation between teacher and student autoregressive sequence models, KL divergence is applied to token-level distributions conditioned on a shared context.
% Let $\pi_\theta(\cdot \mid \mathbf{c})$ and $\pi_{\text{teacher}}(\cdot \mid \mathbf{c})$ denote the student and teacher policies, respectively, defining distributions over the next token given context $\mathbf{c}$.


% Let $\pi_\theta(\cdot \mid \mathbf{c})$ denote the \emph{student} policy and
% $\pi_{\text{teacher}}(\cdot \mid \mathbf{c})$ a fixed \emph{teacher} policy,
% where both define distributions over next tokens conditioned on a context
% $\mathbf{c}$ (prompt plus previously generated tokens).
% The KL divergence between two such conditional
% distributions is defined as
% \begin{equation}
% \mathrm{KL}\!\left(\pi_1 \,\|\, \pi_2\right)
% =
% \mathbb{E}_{x \sim \pi_1(\cdot \mid \mathbf{c})}
% \left[
% \log \frac{\pi_1(x \mid \mathbf{c})}{\pi_2(x \mid \mathbf{c})}
% \right].
% \end{equation}
% KL divergence is non-negative and equals zero if and only if the two
% distributions match almost everywhere.
% Crucially, KL divergence is asymmetric, leading to different optimization
% behavior depending on which policy defines the expectation.

\textbf{Forward KL (Teacher-to-Student).}
The forward KL divergence is defined as an expectation over the teacher's distribution:
\begin{equation}
\label{eq:forward_kl}
\mathrm{KL}
% (\pi_{\tt{te}} \,\|\, \pi_\theta) 
(\pi_{\tt{te}}(\cdot\mid\mathbf{c}) \,\|\, \pi_{\theta}(\cdot\mid\mathbf{c})) = \mathbb{E}_{x \sim \pi_{\tt{te}}(\cdot \mid \mathbf{c})} \left[ \log \frac{\pi_{\tt{te}}(x \mid \mathbf{c})}{\pi_\theta(x \mid \mathbf{c})} \right].
\end{equation}

Minimizing \eqref{eq:forward_kl} is equivalent to standard supervised learning (maximizing likelihood on teacher samples). 
It penalizes the student for assigning low probability to any token the teacher considers likely. 
This induces \emph{mode-covering} behavior, where the student attempts to match the entire support of the teacher, potentially leading to overly diffuse distributions if the student has limited capacity~\cite{minka2005divergence}.

% \paragraph{Forward KL (teacher-to-student).}
% The forward KL divergence from the teacher to the student is given by
% \begin{equation}
% \begin{aligned}
% \mathrm{KL}\!\left(\pi_{\text{teacher}} \,\|\, \pi_\theta\right)
% &=
% \mathbb{E}_{x \sim \pi_{\text{teacher}}(\cdot \mid \mathbf{c})}
% \Big[
% \log \pi_{\text{teacher}}(x \mid \mathbf{c})
% \\
% &\qquad\qquad
% -
% \log \pi_\theta(x \mid \mathbf{c})
% \Big].
% \end{aligned}
% \end{equation}

% Because the expectation is taken over teacher-generated tokens, this objective
% strongly penalizes the student for assigning low probability to tokens that
% the teacher considers likely in a given context.
% As a result, minimizing forward KL encourages the student to allocate
% probability mass wherever the teacher does, regardless of whether the student
% would naturally generate those tokens under its own policy.

% This induces \emph{mode-covering} behavior: the student is pushed to cover
% all modes of the teacher’s distribution, including low-probability or
% ambiguous alternatives.
% While this can improve coverage, it often leads to off-policy training and
% overly diffuse student distributions when capacity is limited.

\textbf{Reverse KL (Student-to-Teacher).}
The reverse KL divergence is defined as an expectation over the student's distribution:
\begin{equation}
\label{eq:reverse_kl}
\mathrm{KL}
% (\pi_\theta \,\|\, \pi_{\tt{te}})
(\pi_{\theta}(\cdot\mid\mathbf{c}) \,\|\, \pi_{\tt{te}}(\cdot\mid\mathbf{c}))= \mathbb{E}_{x \sim \pi_\theta(\cdot \mid \mathbf{c})} \left[ \log \frac{\pi_\theta(x \mid \mathbf{c})}{\pi_{\tt{te}}(x \mid \mathbf{c})} \right].
\end{equation}
Since the expectation is taken over student samples, \eqref{eq:reverse_kl} penalizes generated tokens that the teacher considers unlikely, but ignores teacher modes that the student does not visit. It induces \emph{mode-seeking} behavior, encouraging the student to concentrate probability mass on a single high-likelihood mode of the teacher while ignoring others~\cite{minka2005divergence}.

% \paragraph{Reverse KL (student-to-teacher).}
% The reverse KL divergence from the student to the teacher is defined as
% \begin{equation}
% \begin{aligned}
% \mathrm{KL}\!\left(\pi_\theta \,\|\, \pi_{\text{teacher}}\right)
% &=
% \mathbb{E}_{x \sim \pi_\theta(\cdot \mid \mathbf{c})}
% \Big[
% \log \pi_\theta(x \mid \mathbf{c})
% \\
% &\qquad\qquad
% -
% \log \pi_{\text{teacher}}(x \mid \mathbf{c})
% \Big].
% \end{aligned}
% \end{equation}

% Here, the expectation is taken over student-generated tokens.
% The objective penalizes cases where the student assigns high probability to
% tokens that the teacher deems unlikely, but does not directly penalize tokens
% or modes that the student never samples.

% This leads to \emph{mode-seeking} behavior: minimizing reverse KL
% encourages the student to concentrate probability mass on tokens and modes
% that it already visits and that the teacher assigns high likelihood to,
% while ignoring alternative modes that are unlikely under the student policy.


\subsection{On-Policy Distillation} \label{sec:pre_opd}
% Each term definition

% On-policy distillation (OPD)~\cite{agarwal2024policy} is a post-training method in which a student model learns by matching a teacher’s token-level probability distribution on the student’s own generated sequences. Unlike off-policy distillation, which relies solely on teacher-generated trajectories, OPD trains on on-policy rollouts produced by the student itself. This allows the student to receive token-specific feedback on its own outputs, enabling credit assignment in the exact contexts the student visits. As a result, OPD mitigates compounding errors and exposure bias inherent in purely off-policy imitation.

% Several OPD methods have been proposed that optimize the reverse KL divergence~\cite{gu2023minillm, lu2025onpolicydistillation} to encourage mode-seeking behavior, allowing the student to focus on the teacher’s dominant modes within its limited capacity and thereby improve generation correctness.

% Specifically, recent approach~\cite{lu2025onpolicydistillation} defines a token-level reward based on the log-probability difference between the teacher and the student for each token generated by the student, and optimizes it using policy gradient methods.

On-policy distillation (OPD)~\cite{agarwal2024policy} is a post-training method in which a student model learns by matching a teacher's token-level probability distribution on its own generated sequences. By training on on-policy rollouts rather than teacher-generated trajectories, OPD enables precise credit assignment and mitigates compounding errors inherent in off-policy imitation.

Recent OPD methods optimize the reverse KL divergence~\cite{gu2024minillm, lu2025onpolicydistillation}, encouraging mode-seeking behavior that helps the student focus on the teacher's dominant modes. 
Specifically, given inputs $\mathbf{q}\sim\mathcal{D}$, %and student rollouts , 
we denote $\mathbf{x} = (x_1, x_2, \ldots, x_{|\mathbf{x}|})$ as the student-generated token sequence.
With $\mathbf{c}_t = (\mathbf{q},x_{<t})$ as the context for token $t$, we have $x_t\sim\pi_{\theta}(\cdot\mid\mathbf{c}_t)$. %$\mathbf{x}\sim\pi_{\theta}(\cdot\mid\mathbf{q})$ 
The on-policy reverse-KL objective is then defined as:
\begin{equation}\label{eq:opd_ideal}
\mathbb{E}_{\mathbf{q}\sim\mathcal{D}}\left[
\mathbb{E}_{\mathbf{x}\sim\pi_{\theta}(\cdot\mid\mathbf{q})} \left[
\frac{1}{|\mathbf{x}|}\sum_{t=1}^{|\mathbf{x}|}\mathcal{L}_{t}^{\textrm{RKL}}(\theta;\mathbf{c}_t)
\right]
\right],
\end{equation}
where $\mathcal{L}_{t}^{\textrm{RKL}}(\theta;\mathbf{c}_t)$ is the per-token reverse KL from \eqref{eq:reverse_kl}:
% \begin{equation}
% \mathcal{L}_{t}^{\textrm{RKL}}(\theta;\mathbf{c}_t) = \mathbb{E}_{x_t\sim\pi_{\theta}(\cdot\mid\mathbf{c}_t)}\left[\log\frac{\pi_{\theta}(x_t\mid\mathbf{c}_t)}{\pi_{\tt{te}}(x_t\mid\mathbf{c}_t)}\right].
% \end{equation}
\begin{equation}
\mathcal{L}_{t}^{\textrm{RKL}}(\theta;\mathbf{c}_t) = \mathrm{KL}(\pi_\theta(\cdot\mid\mathbf{c}_t) \,\|\, \pi_{\tt{te}}(\cdot\mid\mathbf{c}_t)).
\end{equation}
In practice, \citet{lu2025onpolicydistillation} use a single-sample Monte‑Carlo estimate of the expectation in \eqref{eq:reverse_kl} and plug it into a policy‑gradient-style update. This is done by defining the token-level reward as the log-probability difference between the teacher and student:
\begin{equation} \label{eq:per_token_reward}
A_t
=
\log \pi_{\tt{te}}(x_t \mid \mathbf{c}_t)
-
\log \pi_{\theta}(x_t \mid \mathbf{c}_t).
\end{equation}
This reward measures how preferred the student-selected token is under the teacher distribution in the same context, assigning positive values when the teacher assigns a higher probability to the token than the student, and negative values otherwise.
The student then essentially optimizes the objective $\max_{\theta}\mathbb{E}_{\pi_{\theta}}\left[\sum_t A_t\right]$.
The result is an effective on-policy distillation method based on policy gradient optimization.

\subsection{OPD with Clipped-Reverse KL} \label{sec:opd_witch_crkl}
To stabilize training, the OPD objective can be implemented with standard PPO-style importance sampling and clipping~\cite{schulman2017proximal}. 
We sample trajectories using a behavior policy $\pi_{\theta_{\text{old}}}$ instantiated from the student policy $\pi_\theta$,
query the teacher for log-probabilities of the sampled student tokens, and
define a per-token advantage $\hat{A}_t=
\log \pi_{\tt{te}}(x_t \mid \mathbf{c}_t)
-
\log \pi_{\theta_{\text{old}}}(x_t \mid \mathbf{c}_t)$, substituting the behavior policy %as the token-level reward 
in~\eqref{eq:per_token_reward}.
% \begin{equation}
% \hat{A}_t
% =
% \log \pi_{\tt{te}}(x_t \mid \mathbf{c}_t)
% -
% \log \pi_{\theta_{\text{old}}}(x_t \mid \mathbf{c}_t).
% \end{equation}
We then update $\pi_\theta$ by minimizing the clipped PPO surrogate over
generated tokens. To simplify notation, we omit  the expectations over $\mathbf{q}$ and $\mathbf{x}$ in \eqref{eq:opd_ideal} and focus on single tokens. The surrogate objective is then 
\begin{equation}
\label{eq:opd}
\mathcal{L}_t^{\text{OPD}}(\theta; \mathbf{c}_t)
=
\mathbb{E}_{x_t\sim\pi_{\theta_{\textrm{old}}}(\cdot\mid\mathbf{c}_t)}
\Big[
%\mathcal{L}^{\text{C-RKL}}_t(\theta;\mathbf{c}_t)
\widetilde{A}_t
\Big],
\end{equation}
where $\widetilde{A}_t$ %$\mathcal{L}^{\text{C-RKL}}_t(\theta)$ 
is the clipped reverse KL loss:
\begin{equation}
\label{eq:clipped}
\begin{aligned}
\widetilde{A}_t%\mathcal{L}^{\text{C-RKL}}_t(\theta;\mathbf{c}_t) 
= \max \big(
 - r_t \hat{A}_t,\, %\\
 - \mathrm{clip}(r_t, 1-\epsilon, 1+\epsilon)\, \hat{A}_t
\big).
\end{aligned}
\end{equation}
%
Here, $r_t = \frac{\pi_\theta(x_t \mid \mathbf{c}_t)}
     {\pi_{\theta_{\text{old}}}(x_t \mid \mathbf{c}_t)}$ corrects for sampling under $\pi_{\theta_{\text{old}}}$, while
clipping limits overly large policy updates. %For notational simplicity, we drop explicit dependence on $\mathbf{c}_t$.
% Unlike reward model based RL, the OPD
% signal is directly tied to matching a trusted teacher distribution and can be
% computed efficiently by evaluating teacher log-probabilities on student rollouts.
